% !TEX root = ../../../main.tex
% 来源：高中物理甲种本·第三册（电磁·光学·近代物理）
% 章节：光的反射和折射 / 透镜成像作图法
% 原始文件：5.tex，行 1421-1447
% 类型：tikz
% ---
    \begin{tikzpicture}[>=latex, scale=.8]
    \draw [pattern=north east lines] (0,2) to [bend left=20] (0,-2) to [bend left=20](0,2);
    \draw [dashdotted] (-5,0)--(5,0);
    \node at (0,0)[fill=white, above]{$O$};
    \draw[thick] (-4,1.5)node[left]{$S$}--(0,1.5)--(2,0)--(4,-1.5)node[right]{$S'$};
    \draw[thick] (-4,1.5) --(0,-1.5)--(4,-1.5);
    \draw[thick] (-4,1.5)--(0,.5)--(4,-1.5);
    \draw[thick] (-4,1.5)--(0,-.5)--(4,-1.5);
    \node at (2,0)[above]{$F'$};
    \node at (-2,0)[below]{$F$};
    \draw (-4,1.5) [fill=black] circle (1.5pt);
    \draw (4,-1.5) [fill=black] circle (1.5pt);
    \draw (0,0) [fill=black] circle (1.5pt);
    \draw (-2,0) [fill=black] circle (1.5pt);
    \draw (2,0) [fill=black] circle (1.5pt);
    
    \draw [->](-3,1.5)--(-1,1.5);
    \draw [->](1,-1.5)--(3,-1.5);
    \draw [->](-2,0)--(-1,-1.5/2);
    \draw [->](2,0)--(3,-1.5/2);

    \draw [->](0,.5)--(2,-.5);
    \draw [>-](-2,1 )--(0,.5);
    \draw [->](0,-.5)--(2,-1);
    \draw [>-](-2,.5 )--(0,-.5);    
    
    \end{tikzpicture}
